% Prior Law References and Statutory Crosswalk (MAIN POINT 3). Built by the
% future pass from the instructions below. Do not process now. This is the
% crosswalk that ties the bill to existing federal and state law and to the prior
% author regulatory adaptations. The dense citations belong in the crosswalk
% tables, not in the prose.

[Write a connecting paragraph for each of the four movements, then carry the
detail in crosswalk tables. Every law must resolve to its bibliography entry so
the rendered reference carries the eCFR, Cornell LII, or legislature URL. Use the
left-aligned ragged-right columns and put the longer parent citation context once
near the top of each table.]

[Movement 1, Federal law crosswalk. Present, in one table, each federal authority
the bill builds on and the precise way the bill extends it without displacing it:
Medicare reasonableness and medical-necessity determinations,
\texttt{usc-1395y} (42 U.S.C. § 1395y); the HIPAA Privacy and Security Rules,
\texttt{cfr-hipaa} (45 CFR Parts 160 and 164); the FDA medical-device
regulations, \texttt{cfr-devices} (21 CFR Parts 860 to 892); and the federal
antidiscrimination guarantees, \texttt{usc-titlevi} (Title VI, 42 U.S.C. § 2000d)
and \texttt{usc-ada} (the Americans with Disabilities Act, 42 U.S.C. § 12101).
State explicitly that these antidiscrimination provisions are codified at 42
U.S.C., and that the bill cites 42 U.S.C. § 2000d and 42 U.S.C. § 12101; an
earlier informal reference to 26 U.S.C. (the Internal Revenue Code) for these
provisions was a mislabel and is corrected here. Add the supporting federal
authorities the bill rests on: the Federal Food, Drug, and Cosmetic Act
(\texttt{usc-fdca}), the Public Health Service Act (\texttt{usc-phsa}), the
electronic-records rule (\texttt{cfr-part11}), and the investigational pathways
21 CFR Part 312 and Part 50 (\texttt{cfr-part312}, \texttt{cfr-part50}). Use the
authority and source lines modelled in
\texttt{template-paper/chunk\_01\_front\_matter.md} and the federal-authority map
in the national-platform \texttt{regulatory\_landscape.tex} and
\texttt{gov\_framework.tex}.]

[Movement 2, State law crosswalk. Present, in one table, the four state statutes
the bill is designed to harmonise with and build on, with one column for the
state requirement and one for how it informs a clause of this bill: New York
Assembly Bill A9149 (\texttt{ny-a9149}), which requires insurers and HMOs to
submit AI algorithms and training data sets to the Department of Financial
Services for certification that they minimize bias on protected characteristics
and adhere to evidence-based clinical guidelines; Texas Senate Bill 1822
(\texttt{tx-sb1822}), noting the governance change from the filed version (submit
the algorithm and training datasets) to the committee substitute (an annual AI
compliance statement, with algorithms required only on a commissioner finding of
noncompliance); California Senate Bill 1120, the Physicians Make Decisions Act
(\texttt{ca-sb1120}), which keeps medical-necessity determinations with a licensed
physician and bars reliance on generalized rather than individual clinical data;
and Florida House Bill 527 with companion Senate Bill 202 (\texttt{fl-hb527}),
which bars AI as the sole basis for a denial and requires independent human
review. Draw the dual state-and-federal jurisdiction treatment from
\texttt{regulatory\_landscape.tex}.]

[Movement 3, Prior author regulatory adaptations. Note that the bill is consistent
with, and can codify portions of, the author's prior Physical AI regulatory
adaptations: the 21 CFR Part 312 adaptation (\texttt{kawchak2026cfr312}, the
template source for this bill), the 21 CFR Part 50 adaptation
(\texttt{kawchak2026cfr50}), the ICH E6(R3) adaptation (\texttt{kawchak2026iche6r3}
and \texttt{iche6r3}), and the national MCP-PAI standard (\texttt{kawchak2026mcp}).
Use the change-summary-table model from
\texttt{template-paper/chunk\_01\_front\_matter.md} to show how a section-by-
section adaptation maps onto a bill, and forward-reference the bill's own
section-by-section analysis.]

[Movement 4, Differing governance and corrections. In a short closing paragraph,
state the governance differences the crosswalk surfaces (the New York submit-and-
certify model versus the Texas attest-unless-noncompliant model versus the
California physician-decision model versus the Florida human-review model), explain
that the bill adopts a verification-before-generation standard that sits upstream
of all of them at the trial-conduct robot-code layer, and confirm the corrected
citations (42 U.S.C. for the antidiscrimination titles; the Texas committee-
substitute as the operative version). Do not duplicate the crosswalk tables here;
reference them.]

% Model federal crosswalk table (citation context at top, abbreviated):
% \begin{table}[h]\centering
% \begin{tabularx}{\textwidth}{L{3.0cm} L{3.6cm} Y}
% \toprule
% \textbf{Authority} & \textbf{Citation} & \textbf{How the bill builds on it} \\
% \midrule
% Medicare & 42 U.S.C. § 1395y & Medical-necessity stays with a licensed human \\
% HIPAA & 45 CFR Parts 160, 164 & Audit-trail privacy and security \\
% Devices & 21 CFR Parts 860--892 & Robot and software device documentation \\
% Title VI & 42 U.S.C. § 2000d & Bias minimization on protected characteristics \\
% ADA & 42 U.S.C. § 12101 & Accessibility and accommodation \\
% \bottomrule
% \end{tabularx}
% \caption{Federal-law crosswalk (corrected antidiscrimination citations at 42 U.S.C.).}
% \end{table}
